\section{Method}
\label{sec:method}

\subsection{Overview}
Each incoming frame $I_t$ is mapped to an embedding $e_t = \phi(I_t) \in
\mathbb{R}^D$ by a frozen backbone $\phi$ (Section~\ref{sec:method-backbone}).
A stack of causal diagonal state-space layers (Section~\ref{sec:method-ssm})
maintains a fixed-size state and produces $\hat{e}_{t+1}$, a prediction of
the \emph{next} embedding, using only $e_1,\dots,e_t$. The anomaly score at
time $t$ is the prediction error
$s_t = \lVert \hat{e}_t - e_t \rVert_2$ (Section~\ref{sec:method-score}),
following the predictive-coding formulation standard in the VAD literature
\cite{liu2018future}, adapted here to operate strictly causally at the
embedding level.

\subsection{Frozen backbone}
\label{sec:method-backbone}
% TODO: report whichever backbone the final numbers use (configs/*.yaml
% `backbone.name`); code supports ImageNet-pretrained ResNet-18
% ($D{=}512$) and DINOv2 ViT-S/14 ($D{=}384$). The backbone is frozen
% throughout training — only the temporal core (below) is optimized.
We deliberately keep $\phi$ simple and frozen: the contribution of this
paper is the temporal core and its analysis, not backbone design, and
freezing $\phi$ keeps the trainable parameter count and compute footprint
small enough to iterate on consumer hardware (Section~\ref{sec:experiments}).

\subsection{Causal diagonal state-space core}
\label{sec:method-ssm}
Each layer maintains a state $s_t \in \mathbb{R}^N$ and updates it as
\begin{align}
  u_t &= W_{\text{in}}\, \mathrm{LN}(x_t) \\
  g_t &= \sigma\big(\mathrm{MLP}([\mathrm{LN}(x_t),\, s_{t-1}])\big) \label{eq:gate}\\
  \bar{a} &= a_{\min} + (a_{\max}-a_{\min})\, \sigma(\theta_a) \label{eq:base-decay}\\
  a_t &= \bar{a} \odot g_t \\
  s_t &= a_t \odot s_{t-1} + u_t \label{eq:recurrence}\\
  y_t &= W_{\text{out}}\, s_t + x_t
\end{align}
where $x_t$ is the layer's input (the backbone embedding $e_t$ for the
first layer, the previous layer's output $y_t$ otherwise), $\mathrm{LN}$
is layer normalization, $\theta_a \in \mathbb{R}^N$ is a learned
per-channel parameter, and $[a_{\min}, a_{\max}] \subset (0,1)$ bounds the
time-invariant \emph{base} decay $\bar{a}$ (Eq.~\ref{eq:base-decay},
following the S4D/S5 initialization convention \cite{gu2022s4,smith2023s5}).
The gate $g_t \in (0,1)^N$ (Eq.~\ref{eq:gate}) is the only
input/state-dependent (``selective'', in Mamba's terminology
\cite{gu2023mamba}) part of the recurrence: it lets the effective decay
$a_t$ drop sharply — fast forgetting, i.e. an implicit state reset — when
the current input is inconsistent with the current state, which is
exactly the situation at an anomaly onset. We call $g_t$ the
\emph{event-boundary gate} for this reason.
Eq.~\ref{eq:recurrence} is a diagonal (real-valued) linear recurrence,
not a full re-implementation of Mamba's selective-scan kernel, which is
CUDA-only; this keeps the entire model runnable on non-CUDA edge hardware
(Section~\ref{sec:experiments}) while retaining $O(N)$ per-step
state size and $O(1)$ per-step compute in the sequence length.

\subsubsection{Strict causality, by construction}
Unlike prior SSM-based VAD methods that process fixed-size clips or
maintain a sliding buffer internally
\cite{vadmamba2025,wavemambaad2025}, Eqs.~\ref{eq:gate}--\ref{eq:recurrence}
depend only on $s_{t-1}$ and $x_t$ — there is no operation in the model
that looks ahead or re-processes past frames. Algorithm~\ref{alg:step}
gives the exact per-frame update used both at training time (called once
per timestep, unrolled over a window) and at deployment (called once per
incoming frame); using the identical routine in both settings is a
deliberate choice, so that reported efficiency is not an artifact of a
training-time-only formulation (see \texttt{src/models/ssm.py}).

\begin{algorithm}[t]
\caption{Streaming per-frame update (one layer)}
\label{alg:step}
\begin{algorithmic}[1]
\Require frame embedding $x_t$, previous state $s_{t-1}$
\State $u_t \gets W_{\text{in}} \, \mathrm{LN}(x_t)$
\State $g_t \gets \sigma(\mathrm{MLP}([\mathrm{LN}(x_t), s_{t-1}]))$
\State $a_t \gets \bar{a} \odot g_t$
\State $s_t \gets a_t \odot s_{t-1} + u_t$
\State $y_t \gets W_{\text{out}} s_t + x_t$
\State \Return $y_t, s_t$
\end{algorithmic}
\end{algorithm}

\subsection{Training objective and anomaly scoring}
\label{sec:method-score}
The stack's final output $y_t$ is passed through a small MLP head to
produce $\hat{e}_{t+1}$, a prediction of the next frame's embedding. The
model is trained self-supervised on normal-only training clips by
minimizing
\begin{equation}
  \mathcal{L} = \tfrac{1}{T-1}\textstyle\sum_{t=1}^{T-1} \lVert \hat{e}_{t+1} - e_{t+1} \rVert_2^2 .
\end{equation}
At inference, the anomaly score at frame $t$ is the prediction error
$s_t = \lVert \hat{e}_t - e_t \rVert_2$ using the prediction made one step
earlier (i.e. using only $e_1,\dots,e_{t-1}$); scores are min-max
normalized per clip before thresholding/AUC computation, following
standard VAD evaluation protocol \cite{liu2018future}.
